\documentclass[10pt,a4paper]{amsart}
\usepackage[margin=19mm]{geometry}
\usepackage{amsmath,amssymb,mathtools}
\usepackage[scr=rsfs,cal=euler]{mathalfa}
\usepackage{xfrac}
\usepackage[unicode,hidelinks,bookmarks=false]{hyperref}
\newcommand{\ie}{i.e.\ }
\newcommand{\cf}{cf.\ }
\newcommand{\resp}{respectively\ }
\newcommand{\Subsection}{\subsection}
\providecommand{\nobreakdash}{\nobreak\hbox{-}}
\setlength{\emergencystretch}{2em}
\multlinegap0pt
\usepackage{bm}
\usepackage{bbm}
\usepackage{dsfont}
\usepackage{xspace}
\usepackage{cancel}
\usepackage{mathtools}
\usepackage{amsthm}
\makeatletter
\@ifundefined{theorem}{\newtheorem{theorem}{Theorem}}{}
\@ifundefined{proposition}{\newtheorem{proposition}[theorem]{Proposition}}{}
\makeatother
\makeatletter
\newcommand{\D}{\mathrm{d}}
\newcommand\E{\mathds{E}}
\newcommand\N{\mathbb{N}}
\newcommand\R{\mathbb{R}}
\newcommand{\V}{\bm{\mathcal{V}}}
\newcommand\Y{\bm{\mathcal{Y}}}
\newcommand{\ceil}[1]{\lceil#1\rceil}
\newcommand\e{\mathord{\mathrm{e}}}
\expandafter\ifx\csname example\endcsname\relax
\newenvironment{example}
    {\pushQED{\qed}\renewcommand{\qedsymbol}{$\triangle$}\examplex}
    {\popQED\endexamplex}
\fi
\renewcommand\ge{\geqslant}
\renewcommand\le{\leqslant}
\newcommand{\mW}{\mathcal{W}}
\expandafter\ifx\csname remark\endcsname\relax
\newenvironment{remark}
    {\pushQED{\qed}\renewcommand{\qedsymbol}{$\Diamond$}\remarkx}
    {\popQED\endremarkx}
\fi
\makeatother
\title[Rates of convergence for multivariate SDEs driven by L\'evy ]{Rates of convergence for multivariate SDEs driven by L\'evy processes in the small-time stable domain of attraction}
\author{Jorge Gonz\'alez C\'azares, David Kramer-Bang}
\date{Source: arXiv:2509.12377v1 (CC BY 4.0)}
\begin{document}
\maketitle

\noindent\emph{Source and licence.} Rates of convergence for multivariate SDEs driven by L\'evy processes in the small-time stable domain of attraction. Version arXiv:2509.12377v1 (CC BY 4.0), distributed under the Creative Commons Attribution 4.0 International licence, as recorded in the arXiv licence metadata for this exact version and shown on the abstract page https://arxiv.org/abs/2509.12377v1. Full rights notice: licenses/arxiv-2509.12377-CC-BY-4.0.txt.\par\smallskip

\noindent\textit{Editorial scope.} Counted unit: the Proposition whose source label is \texttt{prop:simple\_lin\_SDE\_gen\_bound} in the article (source0.tex lines 1395--1407, source label \texttt{prop:simple\_lin\_SDE\_gen\_bound}), delivered together with the article's own complete original proof of it (source0.tex lines 1409--1483, one \texttt{proof} environment of 75 lines). No mathematics is written, repaired or re-derived: statement and proof are the article's own text line for line. Required context retained uncounted: eq:simple\_lin\_SDEs\_prop (lines 1389--1392). Every definition, display and label of those lines is retained unchanged. Omitted from this package and not redistributed: 1--1388, 1393--1394, 1408--1408, 1484--2335. Nothing inside a retained excerpt is omitted or altered: every retained body is delivered exactly as the article's lines are written, so no omission receipt of any kind accompanies this package. Citations: the retained excerpts carry no citation command at all, so no bibliography entry of the article is transcribed and no reference is redistributed. Reference adaptation: none was needed and none was made; every reference inside the delivered unit resolves inside it, so no unresolved reference or citation is produced. Printed slips kept literal: no printed word, symbol, hypothesis or number of the counted statement or of its proof was altered. Layout and class: the article's own class is \texttt{amsart} with packages including geometry, inputenc, mathtools, mathrsfs, multirow, amssymb, makecell, amstext, float, dsfont, amsthm, amsmath, bm, comment; the wrapper is the lane's pinned \texttt{amsart} editorial wrapper at 10pt on A4, which restates the theorem-like environments the retained text needs and adds the article's own macro definitions that the retained text needs (\texttt{\textbackslash D}, \texttt{\textbackslash E}, \texttt{\textbackslash N}, \texttt{\textbackslash R}, \texttt{\textbackslash V}, \texttt{\textbackslash Y}, \texttt{\textbackslash ceil}, \texttt{\textbackslash e}, \texttt{\textbackslash ge}, \texttt{\textbackslash le}, \texttt{\textbackslash mW}), with extra wrapper packages (bm, bbm, dsfont, xspace, cancel, mathtools). The delivered numbering is the unit's own. Assets: no retained span contains a figure environment, a graphics-inclusion command, a PDF-page inclusion, a table or a TikZ picture; no image file is copied into this package, the package declares no asset dependency and no image file is redistributed with it. Licence: version arXiv:2509.12377v1 of this article is distributed under the Creative Commons Attribution 4.0 International licence, as recorded in the arXiv licence metadata for this exact version and shown on the abstract page https://arxiv.org/abs/2509.12377v1. A full rights notice is delivered at licenses/arxiv-2509.12377-CC-BY-4.0.txt. Characters: every retained excerpt is pure ASCII, so the delivered engine, XeLaTeX with its default Latin Modern text font, reports no missing character.
\begin{equation}
\label{eq:simple_lin_SDEs_prop}
\Y(s) = \bm{y} + \int_0^s V(\Y(r-)) \, \mathrm{d}r + \bm{Y}(s),
\end{equation}

\begin{proposition}
\label{prop:simple_lin_SDE_gen_bound} 
Fix $q \in (0,1]$ and, for $i\in\{1,2\}$, let $\bm{Y}_i$ be a L\'evy process with finite $q$-moment and $\Y_i$ be the solution to the SDE in~\eqref{eq:simple_lin_SDEs_prop} driven by $\bm{Y}_i$ with initial value $\bm{y}_i\in\R^d$. Then, for any $T>0$,
\begin{equation*}
\mW_q(\Y_1,\Y_2) 
\le \big(|\Delta\bm{y}|^q +\mW_{q}(\bm{Y}_1,\bm{Y}_2) \big)
\times\begin{dcases}
\e^{KT}, & q=1,\\
 N(N+1) 2^{N-1}, & q \in (0,1),
\end{dcases}
\end{equation*}
where $N \coloneqq \ceil{TK2^{1/q}} \in \N$.
\end{proposition}

\begin{proof}
Part~1. Suppose $q=1$, consider any coupling $(\bm{Y}_1,\bm{Y}_2)$ on $[0,T]$, let $R(t)\coloneqq \E[\|\Delta\Y\|_{[0,t]}]$ and $\V_i\coloneqq V(\Y_i)$. The triangle inequality, the Lipschitz continuity of $V$ and Tonelli's theorem, give
\begin{align*}
R(t) &\le |\Delta\bm{y}|+ \E\|\Delta\bm{Y}\|_{[0,t]}
+\E\bigg\| \int_0^\cdot \Delta\V(s-)\,\D s \bigg\|_{[0,t]} \\
&\le |\Delta\bm{y}|+ \E\|\Delta\bm{Y}\|_{[0,t]} + K \int_0^t \E[\|\Delta\Y\|_{[0,s]}]\, \D s
\le|\Delta\bm{y}| + \E\|\Delta\bm{Y}\|_{[0,T]} + K \int_0^t R(s)\, \D s,
\end{align*} 
for all $t\le T$. Hence, Gr\"onwall's inequality gives
\[
R(t) \le \big(|\Delta\bm{y}|
+ \E\|\Delta\bm{Y}\|_{[0,T]}\big) \e^{Kt},
\quad t\le T.
\]
Setting $t=T$ and taking infimum over all couplings of $(\bm{Y}_1,\bm{Y}_2)$ on $[0,T]$ gives the first claim.

Part~2. Suppose $q \in (0,1)$. A direct application of Gr\"onwall's appears to be impossible, so we will use an iterative bound, resulting in a discrete version of Gr\"onwall's inequality. Again, consider any coupling of $(\bm{Y}_1,\bm{Y}_2)$ on $[0,T]$. By the sub-additivity of $x \mapsto x^q$ for all $x \ge 0$ and the Lipschitz continuity of $V$, 
\begin{align*}
\E[\|\Delta\Y\|_{[0,t]}^q] 
&\le |\Delta\bm{y}|^q
+\E[\|\Delta\bm{Y}\|_{[0,t]}^q] 
+ \E\Bigg[\bigg\|\int_0^\cdot \Delta\V(s-) \,\D s \bigg\|^q_{[0,t]}\Bigg]\\
&\le |\Delta\bm{y}|^q
+\E[\|\Delta\bm{Y}\|_{[0,t]}^q] 
+ (Kt)^q \E\big[\|\Delta\Y(s-)\|^q_{[0,t]}\big], 
\end{align*} 
for all $t\le T$. Hence, for all $t\in (0,K^{-1})$, it holds that
\begin{equation}\label{eq:wasserstein_first_step}
\E[\|\Delta\Y\|_{[0,t]}^q] 
\le (1-(Kt)^q)^{-1} \big(|\Delta\bm{y}|^q 
+ \E\big[\|\Delta\bm{Y}\|_{[0,t]}^q\big]\big).
\end{equation} 
An extension of~\eqref{eq:wasserstein_first_step} for all $t>0$, will follow by iterating~\eqref{eq:wasserstein_first_step} in this case. 

Fix $t^* \in (0,K^{-1})$, consider~\eqref{eq:wasserstein_first_step} applied to $t^*$ and set $C = C_{q,K,t^*} \coloneqq (1-(Kt^*)^q)^{-1} \in (1,\infty)$. The goal is now to iterate the inequality and show that, for all $N \ge 1$,
\begin{equation}\label{eq:induction_N_step}
\E\big[\|\Delta\Y\|_{[0,Nt^*]}^q\big]
\le |\Delta\bm{y}|^q \sum_{k=1}^N C^k + \E\big[\|\Delta\bm{Y}\|_{[0,t^*]}^q\big] \sum_{k=1}^N (N-k+1)C^k.
\end{equation} 
To show this, we do a proof by induction. The base case follows from~\eqref{eq:wasserstein_first_step}. Assume~\eqref{eq:induction_N_step} holds for some $N\ge 1$. It remains to show that~\eqref{eq:induction_N_step} holds for $N+1$. 

Applying the induction hypothesis conditional on the $\sigma$-algebra $\sigma(\bm{Y}_1(s),\bm{Y}_2(s);s\in[0,t^*])$ and integrating, we get
\begin{align*}
\E\big[\|\Delta\Y\|_{[t^*,(N+1)t^*]}^q\big]
&\le\E\big[|\Delta\Y(t^*)|^q\big] \sum_{k=1}^N C^k 
+\E\big[\|\Delta\bm{Y}\|_{[0,t^*]}^q\big]\sum_{k=1}^N (N-k+1)C^k\\
&\le C\big(|\Delta\bm{y}|^q + \E\big[\|\Delta\bm{Y}\|_{[0,t^*]}^q\big]\big) \sum_{k=1}^N C^k 
+\E\big[\|\Delta\bm{Y}\|_{[0,t^*]}^q\big]\sum_{k=1}^N (N-k+1)C^k.
\end{align*}
Hence, we may conclude the induction:
\begin{align*}
\E\big[\|\Delta\Y\|_{[0,(N+1)t^*]}^q\big]
&\le \E\big[\|\Delta\Y\|_{[0,t^*]}^q\big]
+ \E\big[\|\Delta\Y\|_{[t^*,(N+1)t^*]}^q\big]\\
&\le |\Delta\bm{y}|^q \sum_{k=1}^{N+1} C^{k}
+ \E\big[\|\Delta\bm{Y}\|_{[0,t^*]}^q\big] \sum_{k=1}^{N+1} ((N+1)-k+1)C^k.
\end{align*} 

Given $T>0$, set $N \coloneqq \ceil{K2^{1/q}T} \in \N$ and $t^*\coloneqq T/N \in (0,K^{-1})$. Then~\eqref{eq:induction_N_step} gives
\[
\E\big[\|\Delta\Y\|_{[0,T]}^q\big]
\le\big(|\Delta\bm{y}|^q 
    + \E\big[\|\Delta\bm{Y}\|_{[0,T]}^q\big]\big) C^{N} 
\sum_{k=1}^N kC^{1-k},
\enskip\text{where}\enskip
\sum_{k=1}^N kC^{1-k}
\le \frac{N(N+1)}{2},
\]
since $T=Nt^*$ and $t^*\le T$. By construction, $(Kt^*)^q=(KT/N)^q \le 1/2$ and hence $C \in (1,2]$. Thus, 
\begin{equation*}
\E\big[\|\Delta\Y\|_{[0,T]}^q\big]
\le \big(|\Delta\bm{y}|^q +\E\big[\|\Delta\bm{Y}\|_{[0,T]}^q\big]\big) N(N+1) 2^{N-1}. 
\end{equation*}
Taking infima over all couplings of $(\bm{Y}_1,\bm{Y}_2)$, we obtain the claim.
\end{proof}

\end{document}
